\section{Architecture Decision Records}

%%% Texto introductorio a la sección


%%%%%% COMO SOMOS 5, CADA UNO AGREGA UN ADR DE SU FRENTE

% ------------------------- Payments and Refunds -------------------------
\begingroup
\setlength{\emergencystretch}{3em}
\subsection{ADR-PAY-01 (Payments and Refunds): Hold the money at booking and capture it on confirmation}
\begin{description}
  \item[Status] Proposed --- 2026-11-25 (agreed contract with the Reservation
    work stream: \texttt{POST /payments}, \texttt{Payment\allowbreak Succeeded},
    reservation events).
  \item[Context] A booking involves the Reservation Service, Payment
    Coordination and the Ledger, each with its own database, plus an external
    gateway that is slow and cannot join a distributed transaction. The
    Reservation Service holds the dates for 10 minutes and confirms on
    \texttt{Payment\allowbreak Succeeded}. A booking that fails must never leave the
    guest charged (Book, alternative flows; proposed FR11).
  \item[Decision] Payment Coordination \textbf{authorizes} (holds) the money
    when the Reservation Service requests the payment and publishes
    \texttt{Payment\allowbreak Succeeded}. It \textbf{captures} only when it receives
    \texttt{Reservation\allowbreak Confirmed}, and \textbf{voids} the hold on
    \texttt{Reservation\allowbreak Expired} or \texttt{Reservation\allowbreak Cancelled} before
    capture.
  \item[Alternatives] (1) Charge immediately and refund if the reservation
    fails. (2) A distributed transaction across reservation, payment and
    gateway. (3) Payment Coordination calls the Reservation Service to
    confirm (a second orchestrator).
  \item[Rationale] Charging first turns every lost hold into a refund that
    costs fees and takes days to reach the guest; releasing a hold is free and
    immediate. The gateway does not support distributed transactions. Reacting
    to the reservation events keeps the Reservation Service as the only owner
    of the booking state.
  \item[Trade-offs] One more gateway call per booking (capture); an
    authorization expires after some days, so capture must not be delayed;
    states agree only after a few seconds.
  \item[Consequences] Payment Coordination subscribes to the reservation
    events; every operation (authorize, capture, void, refund) is idempotent.
\end{description}
\endgroup

\subsection{ADR-PREV-DOB-BOOK: Prevent double bookings with a database exclusion constraint and expiring holds}
\begin{description}
  \item[Context]
  An accommodation must never be booked twice for overlapping dates under
  concurrent attempts (about 130 booking requests per second at real scale).
  Guests need time to pay without losing the dates, but abandoned payments must
  not block the calendar. Search is derived data and cannot decide
  availability (FR02, FR04).
  \item[Decision]
  Every active period is a calendar entry protected by a PostgreSQL exclusion
  constraint on \texttt{(listing\_id, period)}. Booking inserts a \texttt{HOLD}
  that expires after 10 minutes; expired holds are released by lazy release, a
  scheduled message and a sweeper, all idempotent. State changes are
  conditional updates.
  \item[Alternatives considered]
  \texttt{SELECT ... FOR UPDATE} on the listing row (depends on every code path
  following the protocol); a distributed lock in a cache (can fail
  mid-transaction); a version column only (does not cover overlapping ranges);
  pay first, then reserve (may charge for dates already gone). All rejected.
  \item[Consequences]
  \begin{itemize}
    \item[+] Double bookings are impossible regardless of instance count or
      application bugs; abandoned payments free the dates automatically.
    \item[$-$] Depends on PostgreSQL range types and \texttt{btree\_gist}; a
      hot listing serializes its bookings; holds hide dates for up to 10
      minutes.
  \end{itemize}
\end{description}

% ------------------------- Pricing Engine -------------------------
\subsection{ADR-PRICE-01 (Pricing Engine): Answer the multiplier synchronously from a precomputed zone snapshot}
\begin{description}
  \item[Status] Proposed --- 2026-11-25 (contract with the Reservation work
    stream: \texttt{POST /api/v1/pricing/multiplier} on listing create or
    edit, calendar occupancy API).
  \item[Context] FR03 requires the price of an accommodation to be
    calculated. The Reservation Service needs the multiplier while the host
    saves a listing (timeout of about 500\,ms; a failure returns
    \texttt{502}), and Payments uses it as platform margin. The indicators
    come from three sources with very different rhythms: search events
    (60\,\% of the traffic), the calendar (daily) and an external events API
    (monthly).
  \item[Decision] Precompute demand and context indicators per zone in two
    separate services, deliver them as events to a local snapshot in the
    Pricing Service, and answer each multiplier request with one read of that
    snapshot and the active model version. Model versions are activated only
    after an administrator's approval (human supervision).
  \item[Alternatives] (1) Compute on request, calling the indicator services
    and SerpAPI synchronously. (2) Save the listing without a multiplier and
    set it later through an event. (3) A fixed platform margin with no
    engine.
  \item[Rationale] Alternative 1 chains three services and an external API
    inside a 500\,ms budget, so any of them could block listing saves.
    Alternative 2 leaves listings without a rate and would need the
    Reservation Service to consume and apply a new event. Alternative 3 does
    not satisfy FR03 and ignores demand.
  \item[Trade-offs] Indicators can be up to one day old (demand) or one
    month old (context); one more derived copy (the snapshot) has to be
    monitored; a new model version does not change existing listings until
    they are edited.
  \item[Consequences] The engine is never on the booking or search path; it
    keeps working with the last snapshot if indicator services or the broker
    fail, and it falls back to \texttt{STANDARD} (1.07) for zones without
    data. Every calculation is stored with its inputs, so any multiplier can
    be explained to a host or an auditor.
\end{description}

% ------------------------- Identity, Communication and Internal Operations -------------------------
\begingroup
\setlength{\emergencystretch}{3em}
\subsection{ADR-ID-01 (Identity, Communication and Internal Operations): Deliver notifications asynchronously with at-least-once delivery and idempotent consumers}
\begin{description}
  \item[Status] Proposed --- 2026-11-26.
  \item[Context] The Notification Service must send without duplicates and
    handle failures per channel. Sending a message, creating a case or
    changing an account must not depend on the availability of email, push
    or SMS providers, which are external and can throttle or fail.
  \item[Decision] Business services publish events through a transactional
    outbox. Notification consumes them, applies preferences and templates, and
    sends through per-channel queues with retries, fallback and a dead-letter
    queue. Each notification has the key \texttt{(eventId, userId, channel)},
    inserted in the same local transaction as the outbound task.
  \item[Alternatives] (1) Send inside the originating request. (2) Let each
    service call the providers directly. (3) Assume exactly-once delivery
    from the broker.
  \item[Rationale] Asynchronous delivery isolates provider latency and
    outages and allows backpressure and retries. Exactly-once delivery cannot
    be assumed, so idempotent consumers give an exactly-once \emph{effect}
    on top of at-least-once delivery. The outbox removes the dual-write
    failure between the database and the broker.
  \item[Trade-offs] Delivery is eventual; duplicate suppression and
    delivery-state tracking add complexity; operators need tools for the
    dead-letter queue.
  \item[Consequences] Every consumer must be idempotent; ordering is kept per
    aggregate, not globally; failures are visible per channel; templates and
    preferences become data owned by Notification.
\end{description}

\endgroup
